\documentclass[12pt]{article}

\usepackage[T1]{fontenc}
\usepackage{tgtermes}
\usepackage[scale=0.92]{tgheros}

\usepackage[a4paper, margin=1.25in]{geometry}
\usepackage{setspace}
\onehalfspacing
\setlength{\parindent}{0pt}
\setlength{\parskip}{0.8em}

\usepackage{booktabs}
\usepackage{array}
\usepackage[hidelinks]{hyperref}
\usepackage{titlesec}
\usepackage{enumitem}
\setlist[itemize]{leftmargin=1.6em, itemsep=5pt, topsep=6pt}

\usepackage{newunicodechar}
\newunicodechar{Δ}{\ensuremath{\Delta}}
\newunicodechar{−}{\ensuremath{-}}
\newunicodechar{≤}{\ensuremath{\le}}
\newunicodechar{≥}{\ensuremath{\ge}}
\newunicodechar{≈}{\ensuremath{\approx}}

\titleformat{\section}{\sffamily\large\bfseries}{\thesection}{1em}{}
\titleformat{\subsection}{\sffamily\normalsize\bfseries}{\thesubsection}{1em}{}

\pagestyle{plain}

\title{\sffamily\vspace{-2em} Zero-Shot Topic Labeling: Extensions and Model Comparison\\[0.3em]\large Supervisor Meeting Summary}
\author{\sffamily Daniel Tesfai}
\date{\sffamily \today}

\begin{document}

\maketitle
\vspace{-1.5em}

\section{Starting Point}

\begin{itemize}
    \item Base paper (BTW2025-122): zero-shot topic labeling with Llama 3 8B, 19-topic vocabulary, 2500 papers.
    \item Reported baseline: 50\% success, hallucination 25\% (titles) / 33\% (abstracts).
    \item Two extensions investigated: (1) fixing the matching step, (2) testing newer models and model size.
\end{itemize}

\section{Extension 1: Matching Fix}

\begin{itemize}
    \item Original matching: exact string match only, breaks on quotes or near-miss wording.
    \item Fix: strip punctuation, then word-containment fuzzy match, with the ontology's parent category (\texttt{computer science}) explicitly excluded from auto-resolution.
\end{itemize}

\begin{table}[h]
\centering
\begin{tabular}{lll}
\toprule
Input & Old matching & New matching \\
\midrule
\texttt{"'software engineering'"} & hallucination & software engineering \\
\texttt{computer vision} & hallucination & computer imaging and vision \\
\texttt{computer science} & hallucination & hallucination (correct) \\
\bottomrule
\end{tabular}
\end{table}

At full scale (Qwen3-8B, titles): 103/2500 answers (4.1\%) recovered by the fix, all the literal string \texttt{computer vision}, 83 became successes, 20 became misclassified, 0 stayed unresolved.

\section{Extension 2: Newer Models vs.\ Paper Baseline}

\begin{table}[h]
\centering
\begin{tabular}{lccc}
\toprule
Model & Success & Misclassified & Hallucination \\
\midrule
Paper (Llama 3 8B), titles & 50\% & 25\% & 25\% \\
Paper (Llama 3 8B), abstracts & 50\% & 17\% & 33\% \\
Qwen3-8B, titles & 68.2\% & 29.3\% & 2.5\% \\
Qwen3-8B, abstracts & 70.0\% & 27.2\% & 2.8\% \\
Claude Haiku 4.5, titles & 70.9\% & 27.6\% & 1.5\% \\
Claude Haiku 4.5, abstracts & 74.0\% & 23.3\% & 2.7\% \\
\bottomrule
\end{tabular}
\end{table}

\begin{itemize}
    \item Every newer model beats the paper by a wide margin, driven mainly by hallucination collapsing to under 4\%.
    \item Titles vs.\ abstracts: at native precision, Qwen3-8B and Claude both favor abstracts, matching every Qwen3.5 size tested in Section 4. An earlier GGUF-quantized run of Qwen3-8B had shown the opposite (titles ahead), which in hindsight looks like a precision artifact of that specific backend, not a genuine per-model difference.
\end{itemize}

\section{Extension 3: Does Model Size Matter?}

Same task, same prompt, Qwen3.5 family at five sizes (0.8B to 27B).

\begin{table}[h]
\centering
\begin{tabular}{lcc}
\toprule
Size & Titles & Abstracts \\
\midrule
0.8B & 52.1\% & 53.7\% \\
4B & 75.4\% & 76.7\% \\
9B & 78.6\% & 80.8\% \\
27B & 76.0\% & 77.0\% \\
\bottomrule
\end{tabular}
\end{table}

\begin{itemize}
    \item \textbf{9B achieves the highest accuracy}, not the largest model tested. Scaling helps sharply to 4B, a little more to 9B, then reverses at 27B.
    \item 0.8B already beats the paper's Llama 3 8B baseline (52.1\% vs 50\%), a model ten times its size, from an older generation.
\end{itemize}

\section{Bonus Finding: Quantization Sensitivity Scales With Size}

\begin{itemize}
    \item \textbf{What:} storing and computing a model's weights with fewer bits than the precision it was trained in. Here: 16-bit floating point (bf16) reduced to 4-bit. A lossy compression of the model's parameters, not the model itself.
    \item \textbf{How:} weights are grouped into small blocks, each block gets a scale factor so its original range of values can be represented with only 16 discrete levels (4 bits) instead of the full 16-bit range. Weights are converted back to a working precision (``dequantized'') just before each computation, so this introduces a small, bounded amount of numerical error per weight, not a change to the model's architecture or training.
    \item \textbf{Why:} roughly 4x less memory per parameter. This is what let the 27B checkpoint (needs $\sim$54GB in bf16) fit entirely on a single 24GB GPU, avoiding a much slower fallback where the excess spills into system RAM. It can also increase inference speed, since less data has to move through memory per token generated.
\end{itemize}

\begin{table}[h]
\centering
\begin{tabular}{llccc}
\toprule
Size & Field & bf16 & 4-bit & Difference \\
\midrule
0.8B & titles & 52.1\% & 39.2\% & \textbf{-12.9 pts} \\
0.8B & abstracts & 53.7\% & 27.6\% & \textbf{-26.1 pts} \\
4B & titles & 75.4\% & 74.6\% & -0.8 pts \\
9B & titles & 78.6\% & 80.0\% & +1.4 pts \\
27B & titles & 75.4\% & 76.0\% & +0.6 pts \\
\bottomrule
\end{tabular}
\end{table}

\begin{itemize}
    \item Small models are fragile under 4-bit quantization, 0.8B loses over a quarter of its accuracy edge.
    \item 4B and up are effectively indifferent, gains and losses under 1.5 points.
    \item Found by accident (a script bug applied quantization to every size, not only 27B), but every repeat run came back bit-identical, generation is fully deterministic, so this is a real precision effect, not noise.
\end{itemize}

\section{Full Results Table}

\begin{table}[h]
\centering
\small
\begin{tabular}{lccc}
\toprule
Run & Success & Misclassified & Hallucination \\
\midrule
Paper (Llama 3 8B), titles & 50\% & 25\% & 25\% \\
Paper (Llama 3 8B), abstracts & 50\% & 17\% & 33\% \\
Qwen3.5-0.8B, titles & 52.1\% & 47.3\% & 0.6\% \\
Qwen3.5-0.8B, abstracts & 53.7\% & 43.9\% & 2.4\% \\
Qwen3-8B, titles & 68.2\% & 29.3\% & 2.5\% \\
Qwen3-8B, abstracts & 70.0\% & 27.2\% & 2.8\% \\
Claude Haiku 4.5, titles & 70.9\% & 27.6\% & 1.5\% \\
Claude Haiku 4.5, abstracts & 74.0\% & 23.3\% & 2.7\% \\
Qwen3.5-4B, titles & 75.4\% & 24.0\% & 0.6\% \\
Qwen3.5-4B, abstracts & 76.7\% & 22.3\% & 1.0\% \\
Qwen3.5-9B, titles & 78.6\% & 20.8\% & 0.6\% \\
Qwen3.5-9B, abstracts & 80.8\% & 18.7\% & 0.5\% \\
Qwen3.5-27B, titles (4-bit) & 76.0\% & 24.0\% & 0.04\% \\
Qwen3.5-27B, abstracts (4-bit) & 77.0\% & 22.9\% & 0.1\% \\
\bottomrule
\end{tabular}
\end{table}

\section{Conclusion}

\begin{itemize}
    \item \textbf{Extension 1 (matching fix):} mainly moves the hallucination column, not the model's actual answers. 4.1\% of Qwen3-8B title answers were wrongly counted as hallucination due to fragile exact-string matching; the fix recovered most of them into success. Part of the paper's reported hallucination rate is a measurement artifact, not only a model failure.
    \item \textbf{Extension 2 (newer models):} mainly collapses hallucination, from 25--33\% (paper) to under 4\% for every newer model tested, while success rises 17--30 points and misclassified stays roughly flat or rises slightly. Newer, better instruction-tuned models are far better at staying inside the controlled vocabulary; that alone accounts for most of the improvement over the paper's baseline, not necessarily better topic judgment.
    \item \textbf{Extension 3 (model size):} hallucination stays low and flat across every size tested (0.6\% to 2.4\%, no trend with size), so hallucination is a generation/training effect, not a size effect. What moves with size is the success/misclassified split: 0.8B misclassifies nearly half the time (47.3\%/43.9\%), 9B misclassifies about a fifth of the time. Size affects whether the model picks the \emph{correct} topic, not whether it stays inside the vocabulary, and this effect does not improve indefinitely, it peaks at 9B and reverses at 27B.
    \item \textbf{Taken together:} the three extensions target different parts of the original failure mode. The matching fix and newer models both attack hallucination, one through evaluation robustness, the other through instruction-following, while model size mainly shifts the success/misclassified balance, that is, topic judgment, with diminishing and eventually reversing returns.
\end{itemize}

\section{Open Questions for Discussion}

\begin{itemize}
    \item Which direction next: cost optimization and multi-label evaluation, or few-shot prompting?
    \item Worth running Qwen3.5-2B to fill the one gap in the size sweep?
    \item How to frame the 9B-highest-accuracy finding in the thesis narrative?
\end{itemize}

\end{document}
